%! Author = lili
%! Date = 5/7/26

\begin{table}[H]
    \centering
    \begin{tabular}{p{6cm} p{6cm}}
        \textbf{Prokaryoten} & \textbf{Eukaryoten} \\
        \hline

        kernlose Zellen &
        kernhaltige Zellen \\

        evolutionär älter &
        Endosymbionten-Theorie (Organellen, ausgeprägte Kompartimentierung) \\

        Genorganisation in Operons (polycistronische mRNA) &
        Genomorganisation in Genen (i.d.R.\ monocistronische mRNA) \\

        Genexpression: Promotor/Operator &
        komplexe Genexpression (RNA Pol-II) \\

        wenig bis keine Introns &
        Gene mit Introns als Regel \\

        Translation am nascenten Transkript &
        Polyadenylierung, Cap-Struktur \\

        „einfache“ Zellteilung &
        Trennung von Transkription und Translation in Kern und Cytoplasma \\

        &
        Mitose, Meiose \\
    \end{tabular}
    \caption{Prokaryoten vs. Eukaryoten}\label{tab:prokaryoten-vs-eukaryoten}
\end{table}

\section{Gene}
Gene liegen auf Chromosomen und sind linear angeordnet.
\paragraph{Was ist ein Gen?}
Naive Annahme:

\begin{center}
    Ein Gen codiert für ein Protein.
\end{center}
Das ist jedoch nicht ganz ausreichend.
Differenzierter könnte man sagen:
\begin{center}
    Ein Gen umfasst den gesamten Bereich auf der DNA, der zur Erstellung
    einer (m)RNA notwendig ist, einschliesslich der regulatorischen Regionen.
\end{center}
Das passt ``so halb'' für Prokaryoten, bei denen Gene oft hinter Operons liegen und daher durch das Operon an einem
Stück transkribiert wird.
Es entsteht eine mRNA, aber diese kann immer noch für mehrere Proteine kodieren.

Für Eukaryoten passt es nicht.
Denn es bezieht nicht ein, dass bei Eukaryoten Gene ``ineinander'' liegen können.
Ein DNA-Bereich kann bei einem Gen ein Intron sein und bei einem anderen Gen ein kodierender Bereich.
Aufgrund der Leseraster kann die DNA-Sequenz auf verschiedene Arten gelesen werden.
Außderdem können Exons beim Spleißen verschieden zusammen gefügt werden.

\frage{wasistgen}

\begin{bem}
    Bei Bakterien und Archaeen sind
    Genzahl und Genomgröße proportional.
    Bei Eukaryoten:
    Keine Korrelation der Genzahl mit Genomgröße oder Komplexität des Organismus.
\end{bem}

Wie findet man sich in so viel DNA bzw. Genen zurecht?

\begin{itemize}
    \item Restriktionskartierung (Beispiel: DNA Tumorviren, etwa Adenovirus-2)
    \item BAC-Contigs (Beispiel: Physikalische Karte von Arabidopsis thaliana)
\end{itemize}

% TODO
\frage{wiekarte}